% Level 2, Week 12 -- Showcase and Celebration (student workbook)
\chapterweek{12}{Showcase and Celebration}{Share what you built, explain how it works, and see how far you have come.}

\section{What you will do today}
\begin{itemize}
  \item Present your project in about a minute: what it is, how one part works, and what you would add next.
  \item Play everyone's projects in a gallery walk and leave kind, specific feedback.
  \item Show what you learned on the exit assessment (it is not a test you can fail: it is your stats screen).
  \item Look back at your Level 2 journey, and peek at Level 3.
\end{itemize}

\section{The big idea: look how far you have come}
Twelve weeks ago you were sorting values into types. Now you have built a whole project with
functions, lists, dictionaries, loops, files and tests. Here is the road you travelled:

\begin{center}
\begin{tikzpicture}[font=\footnotesize, align=center,
  stop/.style={circle, draw=kcblack, thick, minimum size=0.95cm, font=\small\bfseries},
  lbl/.style={text width=2.2cm, font=\footnotesize},
  arr/.style={-Stealth, thick, kcgrey}]
  \foreach \i/\c/\t in {1/kcorange/Types and f-strings, 2/kcsky/Functions that return, 3/kcgreen/Lists, 4/kcyellow/While and game loops, 5/kcblue/Dictionaries and codes, 6/kcpurple/Debugging detectives} {
    \node[stop, fill=\c!35] (s\i) at ({(\i-1)*2.55}, 0) {\i};
    \node[lbl, below=0.1cm of s\i] {\t};
  }
  \foreach \i/\c/\t in {7/kcorange/Files: saving stuff, 8/kcsky/Testing: prove it works, 9/kcgreen/Event-driven games, 10/kcyellow/Project Studio 1, 11/kcblue/Project Studio 2, 12/kcpurple/Showcase!} {
    \node[stop, fill=\c!35] (s\i) at ({(12-\i)*2.55}, -2.6) {\i};
    \node[lbl, below=0.1cm of s\i] {\t};
  }
  \foreach \a/\b in {1/2,2/3,3/4,4/5,5/6,7/8,8/9,9/10,10/11,11/12} { \draw[arr] (s\a) -- (s\b); }
  \draw[arr] (s6.east) to[out=0, in=0, looseness=1.6] (s7.east);
  \node[left=0.15cm of s12, font=\small\bfseries] {★ today};
\end{tikzpicture}
\end{center}

\section{How to explain your project}
A great showcase pitch is short. Answer four questions, in this order:
\begin{enumerate}
  \item \textbf{What is it?} \enquote{This is Comet Catch. You catch comets with a rocket.}
  \item \textbf{How does one part work?} Pick one function, the loop, the list, the dictionary or the file. \enquote{Our \texttt{next\_x} function gets the position and the direction and returns the new position, but never off the screen. We have a test for that.}
  \item \textbf{One design decision.} Why did you choose something? \enquote{The comets are circles and the rocks are squares, so you can tell them apart even without colour.}
  \item \textbf{What next?} \enquote{Next we would add a slow mode for beginners.}
\end{enumerate}
If something is not finished, that is normal. Show what works, then say what you would add next.

\begin{exercise}[My one-minute pitch]
Fill in your presentation card, then practise once with your partner.

What is it? \answerlines{1}
How does one part work? (name a function, the loop, a list, a dictionary or the file) \answerlines{2}
One design decision and why: \answerlines{1}
What would we add next? \answerlines{1}
\end{exercise}

\begin{exercise}[Gallery walk feedback]
Visit every project. Play it for 2--3 minutes, then leave a card: \textbf{two stars} (two specific
things that are good) and \textbf{a wish} (one kind idea). Specific is kind: \enquote{the score
text is big and easy to read} helps more than \enquote{cool}.

\begin{tcolorbox}[colback=white,colframe=kcpurple!80!black,boxrule=0.8pt,arc=2pt,title={\small\bfseries Feedback card for: \rule{4cm}{0.4pt}},coltitle=black,colbacktitle=kcpurple!25]
★ \answerlines{1}
★ \answerlines{1}
Wish: \answerlines{1}
\end{tcolorbox}
\begin{tcolorbox}[colback=white,colframe=kcpurple!80!black,boxrule=0.8pt,arc=2pt,title={\small\bfseries Feedback card for: \rule{4cm}{0.4pt}},coltitle=black,colbacktitle=kcpurple!25]
★ \answerlines{1}
★ \answerlines{1}
Wish: \answerlines{1}
\end{tcolorbox}
\end{exercise}

\begin{exercise}[My Level 2 skills check]
Be honest: \enquote{not yet} just means \enquote{not yet}. Write a tick in one column for each line.

\renewcommand{\arraystretch}{1.45}
\begin{tabularx}{\linewidth}{|X|p{1.7cm}|p{1.9cm}|p{1.9cm}|}
\hline
\textbf{I can…} & \textbf{Not yet} & \textbf{Getting there} & \textbf{I could teach it} \\ \hline
use int, float, str and bool, and build messages with f-strings & & & \\ \hline
write a function with parameters that returns a value & & & \\ \hline
keep many things in a list and loop over it & & & \\ \hline
write a \texttt{while} loop that stops when it should & & & \\ \hline
look things up in a dictionary & & & \\ \hline
read a traceback and find the bug & & & \\ \hline
save data to a file and load it again & & & \\ \hline
write a pytest test and read PASSED / FAILED & & & \\ \hline
make a turtle game react to keys and a timer & & & \\ \hline
plan with pseudocode and share the work fairly & & & \\ \hline
\end{tabularx}
\end{exercise}

\begin{exercise}[Looking ahead to Level 3]
In Level 3 you will save your code with \textbf{git save points} (like save points in a video game),
design your own \textbf{classes} (blueprints for game characters and objects), and build bigger
projects. The example project is \textbf{Maze Muncher}: a maze game where a ghost hunts you down with a
clever search algorithm.

Three things I want to add to my project in a \enquote{version 2}:
\answerlines{3}
\end{exercise}

\begin{safety}[Sharing safely]
Show your project to family and friends you know. Use your nickname on the screen and in high-score
files. If someone takes a photo or video, film the screen, not faces. Never upload or post your
project without deciding together with a parent.
\end{safety}

\begin{vocab}
  \vocabword{showcase}{an event where makers show their work and explain it.}
  \vocabword{pitch}{a short talk that explains what your project is and why it is interesting.}
  \vocabword{design decision}{a choice you made on purpose, and the reason for it.}
  \vocabword{gallery walk}{visiting projects one by one, like in a museum, and leaving feedback.}
  \vocabword{version}{one stage of a project; version 2 is the next, improved one.}
\end{vocab}

\begin{reflection}
What are you proudest of in your project?
\answerlines{2}
What did you learn about working in a team?
\answerlines{2}
What do you want to build next?
\answerlines{2}
\end{reflection}
